\subsection{A special PD conormal of infinite projective dimension}
\label{arch136:special-conormal-n3}

\paragraph{ARCH-136-SPECIAL-CONORMAL-N3: statement and scope.}
Put
\[
 R=\mathbf Z_{(2)},\qquad S=\Gamma_R(T),\qquad
 I=\bigoplus_{w\geq3}R\,T^{[w]},\qquad A=S/I.
\]
Write \(I^{[2]}\) for the ideal generated by \(I^2\) and all
\(\gamma_r(a)\), \(a\in I,\ r\geq2\), and put
\[
                       P=I/I^{[2]}.
\]
This is the special PD conormal: in the Beck conormal \(I/I^2\),
it is the quotient by the \(A\)-span of the prime-power operator
images. Then
\[
                         \operatorname{pd}_A(P)=\infty.
                         \tag{136.SC1}
\]
The statement concerns the underived conormal of one coefficient
tail. It excludes a universal assertion that all such special
conormals admit bounded projective resolutions. It neither computes
the full derived special cotangent complex nor obstructs arbitrary
nonlinear or unbounded PD repairs, and it does not decide integral
connectedness.

The link to the sheets is the simultaneous divided-power-envelope
repair problem in cote~136, physical pp.~117,119,121, footers
116,118,120, supplied scans
\texttt{ARCHIVE\_136\_scan117.tex},
\texttt{ARCHIVE\_136\_scan119.tex},
\texttt{ARCHIVE\_136\_scan121.tex}. This is a new calculation in
the investigation, without an assertion of historical priority.
It uses no input A, M, or outside preprint.

\paragraph{The exact module presentation.}
Let
\[
 x=T,\qquad y=T^{[2]},\qquad
 e=xy=3T^{[3]},\qquad f=3T^{[4]}.
\]
Then
\[
 A=R[x,y]/(x^2-2y,\ xy,\ y^2),
 \tag{136.SC2}
\]
and \(A\) is \(R\)-free with basis \(1,x,y\). The special conormal
has the presentation
\[
 P=
 \frac{Ae\oplus Af}
 {(xe-4f,\ ye-2xf,\ 2yf)}.
 \tag{136.SC3}
\]
The generator weights are \(3,4\). In particular, its underlying
\(R\)-module is
\[
             R e\ \oplus R f\ \oplus R(xf)\
                         \oplus (R/2)(yf).
 \tag{136.SC4}
\]
The summands have weights \(3,4,5,6\), respectively.

\emph{Proof of the presentation.}
The divided-power multiplication formula gives \(x^2=2y\).
All weights at least three vanish in \(A\), proving (136.SC2).
In \(P\),
\[
 xe=x^2y=2y^2=12T^{[4]}=4f,\qquad
 ye=xy^2=6xT^{[4]}=2xf.
\]
Moreover
\[
 \gamma_2(e)=9\gamma_2(T^{[3]})
            =90T^{[6]}=2yf,
\]
which vanishes in \(P\). Thus all three displayed relations hold.

Here is a weightwise completeness check. In weights \(3,4,5\)
there are no relations from \(I^2\) or higher divided powers of
elements of \(I\); these components are free of rank one over \(R\).
In weight six, \(I^2\) is generated by
\((T^{[3]})^2=20T^{[6]}\), while
\(\gamma_2(T^{[3]})=10T^{[6]}\).
The quotient is therefore \(R/2\), and \(yf=45T^{[6]}\) is a
generator. In weight seven, the product
\(ef=315T^{[7]}\) has a unit coefficient and kills that component.

All components of weight at least eight also vanish. Indeed
\[
                       \gamma_2(f)=315T^{[8]}
\]
puts \(T^{[8]}\) in \(I^{[2]}\). Over \(R\), the weight tail starting
at eight is the PD ideal generated by \(T^{[8]}\).
For a direct verification, write \(w=8a+b\), \(0\leq b<8\).
The coefficient of
\(\gamma_a(T^{[8]})T^{[b]}\) in \(T^{[w]}\) is a \(2\)-adic unit:
Legendre's formula makes
\((8a)!/(a!(8!)^a)\) a unit, and adding \(b\) to \(8a\) involves
no binary carry. The ideal \(I^{[2]}\) is a PD ideal. To check
this without an extra assumption, divided powers of its generators
stay in it: powers of products of two elements of \(I\) remain in
\(I^2\), while
\(\gamma_s(\gamma_r(a))\) is an integer multiple of
\(\gamma_{sr}(a)\); the sum and scalar formulas then apply.
It follows that the entire tail starting at eight lies in
\(I^{[2]}\).

Finally the abstract module on the right of (136.SC3) reduces,
using \(xe=4f\), \(ye=2xf\), to the four summands in (136.SC4),
with the single remaining relation \(2yf=0\).
For example the action on this proposed \(R\)-module is
\[
 \begin{array}{c|cccc}
       &e&f&xf&yf\\ \hline
 x     &4f&xf&2yf&0\\
 y     &2xf&yf&0&0 ,
 \end{array}
\]
and it satisfies \(x^2=2y,\ xy=y^2=0\).
Hence these are all the abstract module relations. Its four
weight components map isomorphically to those just calculated
in \(P\), proving (136.SC3)--(136.SC4).
\hfill\(\square\)

\paragraph{Elementary proof of infinite projective dimension.}
The ring \(A\) is noetherian and local, with maximal ideal
\(\mathfrak m=(2,x,y)\). Let
\[
 \bar A=A/2A
       =\mathbf F_2[x,y]/(x,y)^2,\qquad
 \mathfrak r=(x,y)\subset\bar A.
 \tag{136.SC5}
\]
Its radical is nonzero and satisfies \(\mathfrak r^2=0\).
Suppose \(P\) had finite projective dimension. A finite minimal
free resolution can then be chosen: the module is finitely
generated and \(A\) is noetherian local. Let its last nonzero
free term be \(F_\ell\).

Every \(A\)-free module is \(R\)-free, so the same resolution,
viewed over \(R\), computes
\(\operatorname{Tor}^R_*(P,\mathbf F_2)\).
Formula (136.SC4), or the usual two-term resolution of \(R/2\),
shows that these groups vanish in degrees at least two.
After reduction modulo two the differentials of the minimal
resolution have entries in \(\mathfrak r\). If \(\ell\geq2\),
the nonzero subspace \(\mathfrak r(F_\ell/2F_\ell)\) is killed
by its differential, because \(\mathfrak r^2=0\), and there
is no incoming differential at this last term. Thus reduced
homology in degree \(\ell\) is nonzero, a contradiction.
Consequently a finite resolution would have length at most one.

The module is not \(A\)-projective, since it has nonzero
\(R\)-torsion whereas any \(A\)-projective is \(R\)-flat.
It needs precisely two generators over the local ring \(A\):
both \(e\) and \(f\) survive independently modulo
\(\mathfrak mP\). Its \(R\)-rank is three, while \(A\) has
\(R\)-rank three. A length-one minimal resolution must therefore
be
\[
                 0\longrightarrow A
                   \longrightarrow A^2\longrightarrow P
                   \longrightarrow0.
                 \tag{136.SC6}
\]
On the other hand (136.SC3) gives
\[
                     P/2P\cong\mathbf F_2 e\oplus\bar A f.
                     \tag{136.SC7}
\]
This is four-dimensional over \(\mathbf F_2\); since
\(\bar A^2\) is six-dimensional, the image of the reduced
first map in (136.SC6) would have dimension two.

No \(\bar A\)-linear map \(\bar A\to\bar A^2\) has such an image.
If its generating vector has a unit coordinate, projection onto
that coordinate shows it is injective, with three-dimensional
image. Otherwise both coordinates lie in \(\mathfrak r\), and
\(\mathfrak r^2=0\) shows that its image has dimension at most one.
This contradiction proves (136.SC1). No use of
Auslander--Buchsbaum is needed.
\hfill\(\square\)

\paragraph{Comparison with the large-prime calculation.}
For a fixed cutoff \(N\) and \(p>2N\), the separate result
ARCH-136-TAIL-CONORMAL gives a special conormal of projective
dimension one, despite infinite Beck conormal projective
dimension. The present \(N=3,p=2\) example shows that the special
quotient need not have that small projective dimension at an
exceptional prime. It does not contradict its uniform
length-one resolution after inverting \((2N)!\).

\paragraph{Source and checks.}
Only the identification of special conormal versus the Beck
operator quotient uses terminology from A.~Kmail, J.~Kozak and
H.~Miller, \emph{Divided powers and K\"ahler differentials},
\emph{Communications in Algebra} \textbf{54}(5) (2026), 2285--2293,
DOI \texttt{10.1080/00927872.2025.2578689}.
Proposition~4.2 and its \(p\)-local specialization are at printed
pp.~5--6 of the author version
\url{https://math.mit.edu/~hrm/papers/dp5.pdf}.
The explicit weight calculation and homological proof above are
given in full and do not depend on a theorem about a derived
cotangent complex.

The complete proof was checked by the author, root, and independently
by \texttt{/root/cote123} during the current run on
2026-11-22 (GPT-6 Astra Pro; effort not exposed).
All three readers checked the saved module, and the independent
reviewer also checked the earlier communicated proof. That reviewer
supplied the elementary final-resolution
argument replacing an initial appeal to Auslander--Buchsbaum.
The source statement was read directly in the primary author PDF.
These checks are whole-proof readings, not formal verification
or independent rediscovery.
